\SourceLeaf{26}
er diesen zweiten Teil behandelte, ist weitaus verwickelter. Zum Teil stützt sie sich, wie die Methode, die der bedeutende \textsc{Laplace} für äußere Punkte verwendet hatte, auf die Betrachtung unendlicher Reihen, die nicht immer konvergieren. Das hätte sich gewiss vermeiden lassen. Bei näherer Prüfung wird sich im Übrigen zeigen, dass diese Lösung des angesehenen \textsc{Ivory}, die bei flüchtiger Betrachtung eine gewisse Ähnlichkeit mit unserer zu haben scheinen könnte, auf \textit{völlig anderen} Grundsätzen beruht. Beide Lösungen haben kaum etwas gemeinsam außer der Verwendung der von uns mit $p,q$ bezeichneten Unbestimmten.
